\subsection{Monte Carlo Simulation}
\label{subsec:monte_carlo}

The Monte Carlo stage turns the fixed demand and the candidate detours into
repeated SUMO observations. In this stage, the demand stays fixed while the
route assignment changes. The scenario generator gives the simulation an
unrouted \texttt{trips.xml} file and each Monte Carlo iteration routes that
same trip set again with \texttt{duarouter}, a different seed and randomized
edge weights.

Here, a Monte Carlo run means one repetition of the same experiment under a
different routing random seed. The baseline in each run is the route
realization without applying any detour plan or closure treatment. Candidate
plans are then applied to a copy of that same baseline route file. This paired
design means that a plan is compared against the exact route-choice
realization from which it was produced, rather than against an average from
other seeds.

\subsubsection{Experimental design}
\label{sssec:mc_design}

The experiment uses a \textbf{paired stochastic-routing design}. The trip
population is fixed, while the route realization changes from run to run. Let
$\mathcal{T}$ be the \texttt{trips.xml} file produced by the scenario
generator. For Monte Carlo iteration $r$, the simulation first computes the
baseline route file
\begin{equation}
  \mathcal{R}_0^{(r)}
  =
  \textsc{DuaRouter}\bigl(\mathcal{T},\,G,\,r,\,\alpha\bigr),
  \label{eq:duarouter_run}
\end{equation}
where $G$ is the SUMO network and $\alpha$ is the upper bound passed to
\path|--weights.random-factor|. With this option, \texttt{duarouter}
perturbs edge costs by drawing multipliers in $[1,\alpha]$. This does not
change the trip origin or destination, it changes only how attractive each
road segment appears during route computation. In human terms, this represents
a driver who still wants to reach the same destination but may slightly adapt
the chosen path after observing or anticipating traffic on particular streets.
Different iterations can therefore produce different paths for the same trip
endpoints.
Figure~\ref{fig:weights-random-factor} illustrates this mechanism on a small
two-route graph.

\begin{figure}[H]
\centering
\begin{tikzpicture}[
  vx/.style={draw, circle, fill=blue!10, thick, minimum size=7mm,
             font=\small\bfseries, inner sep=0pt},
  sel/.style={draw, very thick, blue!70!black,
              -{Stealth[length=6pt,width=5pt]}},
  unsel/.style={draw, thick, gray!60,
                -{Stealth[length=6pt,width=5pt]}},
  wlbl/.style={font=\footnotesize, fill=white, inner sep=1pt},
  ttl/.style={font=\small\itshape}
]
%--- LEFT: no perturbation ------------------------------------------------
\begin{scope}
  \node[vx] (A)  at (0, 0) {$A$};
  \node[vx] (B)  at (2, 1) {$B$};
  \node[vx] (C)  at (2,-1) {$C$};
  \node[vx] (D)  at (4, 0) {$D$};

  \draw[sel]   (A) -- node[wlbl,above left]  {2} (B);
  \draw[sel]   (B) -- node[wlbl,above right] {2} (D);
  \draw[unsel] (A) -- node[wlbl,below left]  {3} (C);
  \draw[unsel] (C) -- node[wlbl,below right] {3} (D);

  \node[ttl] at (2,-2.1) {No perturbation ($\times\,1$)};
\end{scope}
%--- RIGHT: with perturbation ---------------------------------------------
\begin{scope}[xshift=6.4cm]
  \node[vx] (A2) at (0, 0) {$A$};
  \node[vx] (B2) at (2, 1) {$B$};
  \node[vx] (C2) at (2,-1) {$C$};
  \node[vx] (D2) at (4, 0) {$D$};

  \draw[unsel] (A2) -- node[wlbl,above left]  {\textcolor{red!70!black}{3.8}} (B2);
  \draw[unsel] (B2) -- node[wlbl,above right] {\textcolor{red!70!black}{3.6}} (D2);
  \draw[sel]   (A2) -- node[wlbl,below left]  {\textcolor{red!70!black}{3.3}} (C2);
  \draw[sel]   (C2) -- node[wlbl,below right] {\textcolor{red!70!black}{3.0}} (D2);

  \node[ttl] at (2,-2.1) {With perturbation ($\alpha = 2$)};
\end{scope}
\end{tikzpicture}
\caption{Effect of \texttt{--weights.random-factor} on route choice. Both
  graphs share the same network and origin--destination pair $A\!\to\!D$.
  \textbf{Left}: nominal edge costs, the router selects $A\!\to\!B\!\to\!D$
  (cost~4) over $A\!\to\!C\!\to\!D$ (cost~6). \textbf{Right}: each cost is
  multiplied by a random factor drawn from $[1,\alpha]$, the top path becomes
  costlier (7.4) than the bottom path (6.3), so the router now chooses
  $A\!\to\!C\!\to\!D$. Across Monte Carlo iterations, different draws of
  these multipliers yield different route realizations for the same trip set.}
\label{fig:weights-random-factor}
\end{figure}

For every candidate deviation plan $\mathcal{D}_k$, the deterministic rewriter
is then applied to that same baseline route file. The baseline and all plan conditions are paired within the same
run: they share the same trip set, the same route realization before applying
the closure treatment and the same SUMO seed.

Meanwhile all available metrics are collected and saved for later evaluation. For each iteration $r$, the raw time-loss difference
\begin{equation}
  \delta_k^{(r)} \;=\; D_k^{(r)} \;-\; D_0^{(r)}
  \label{eq:paired_delta}
\end{equation}
measures the effect of plan $\mathcal{D}_k$ relative to the no-closure
baseline under the same route-choice and SUMO random seed. The estimator
\begin{equation}
  \bar{\delta}_k = \frac{1}{R}\sum_{r=1}^{R}\delta_k^{(r)}
\end{equation}
compares plans after much of the run-level noise shared by the baseline and
plan conditions has been removed.

\subsubsection{Per-run routing and detour application}
\label{sssec:mc_precomp}

There is no plan-specific route pre-computation outside the Monte Carlo loop.
The only fixed demand input is \texttt{trips.xml}. In each iteration,
\texttt{duarouter} writes \texttt{run\_\#/routes.xml}. This is the baseline
route realization for that seed. The route rewriter then creates the
\texttt{routes\_\{plan\_id\}.xml} files in the same run directory.

After each plan file is written, the code checks that the closed edge no
longer appears in any rewritten route. It also compares the baseline and plan
files vehicle by vehicle to count how many routes changed. These statistics
are stored per run, since the affected vehicles can change when
\texttt{duarouter} produces a different baseline route set.

\subsubsection{Iteration structure}
\label{sssec:mc_iteration}

Each Monte Carlo iteration consists of routing, deterministic route rewriting,
SUMO execution and metric extraction. Algorithm~\ref{alg:mc_iter} states the
procedure formally.

\begin{algorithm}[htbp]
\caption{\textsc{MonteCarloIteration}}\label{alg:mc_iter}
\begin{algorithmic}[1]
\Require network file $G$,
         trip file $\mathcal{T}$,
         closed edge $e_c$,
         candidate plans $\{\mathcal{D}_1, \ldots, \mathcal{D}_K\}$,
         edge-endpoint map $\mathcal{E}$,
         run index $r$,
         simulation window $[b,\, e]$,
         random factor $\alpha$
\Ensure  per-condition metric dictionary for this iteration
\State $\mathcal{R}_0^{(r)} \leftarrow
       \textsc{DuaRouter}(\mathcal{T},G,r,\alpha)$
       \Comment{baseline route realization}
\State $M_0 \leftarrow \textsc{RunSUMO}(G,\mathcal{R}_0^{(r)},r,b,e)$
\State $\text{results}[\,\text{baseline}\,] \;\leftarrow\; M_0$
\For{$k \leftarrow 1$ \textbf{to} $K$}
  \State $\mathcal{R}_k^{(r)} \leftarrow
         \textsc{ApplyDetourToRoutes}(\mathcal{R}_0^{(r)},\mathcal{D}_k,e_c,\mathcal{E})$
  \State validate that $\mathcal{R}_k^{(r)}$ contains no closed edge $e_c$
  \State $s_k^{(r)} \leftarrow
         \textsc{CountRerouted}(\mathcal{R}_0^{(r)},\mathcal{R}_k^{(r)})$
  \State $M_k \leftarrow \textsc{RunSUMO}(G,\mathcal{R}_k^{(r)},r,b,e)$
  \State $\text{results}[\mathcal{D}_k] \;\leftarrow\; M_k \;\cup\; s_k^{(r)}$
\EndFor
\State \Return results
\end{algorithmic}
\end{algorithm}

\textsc{DuaRouter} invokes \texttt{duarouter} with the fixed trip file,
\texttt{--weights.random-factor} and the run seed. \textsc{RunSUMO} writes a
minimal SUMO configuration file that specifies the network, the route file,
the simulation window $[b, e]$ and the same seed $r$, invokes \texttt{sumo}
as a subprocess and parses the resulting \texttt{tripinfo.xml} output into
the metrics of Section~\ref{sssec:mc_metrics}. Each condition runs in its own
subdirectory (\texttt{run\_r/baseline/}, \texttt{run\_r/plan\_k/}), keeping
outputs isolated.

Figure~\ref{fig:mc-workflow} summarizes the same iteration as a data flow from
fixed inputs to per-run baseline and plan outputs.

\begin{figure}[H]
\centering
\begin{tikzpicture}[
  data/.style={draw, dashed, rounded corners=4pt, fill=orange!7, thick,
    minimum width=9.4cm, text width=8.8cm, minimum height=0.70cm,
    align=center, font=\small},
  stage/.style={draw, rounded corners=4pt, fill=blue!7, thick,
    minimum width=9.4cm, text width=8.8cm, minimum height=0.78cm,
    align=center, font=\small},
  result/.style={draw, rounded corners=4pt, fill=teal!10, thick,
    minimum width=9.4cm, text width=8.8cm, minimum height=0.70cm,
    align=center, font=\small},
  arr/.style={-{Stealth[length=6pt]}, thick},
  runlabel/.style={font=\small\bfseries},
  node distance=0.52cm
]
  \node[data] (inputs)
    {fixed inputs: network $G$, trip demand $\mathcal{T}$ and candidate plans $\{\mathcal{D}_k\}$};

  \node[stage, below=1.25cm of inputs] (route)
    {1. \textsc{DuaRouter}: create one baseline route realization
     $\mathcal{R}_0^{(r)}$};
  \node[stage, below=of route] (rewrite)
    {2. route rewriting: apply every plan to the same $\mathcal{R}_0^{(r)}$};
  \node[stage, below=of rewrite] (sim)
    {3. SUMO: simulate baseline and plan routes with the same seed $r$};

  \begin{scope}[on background layer]
    \node[draw, rounded corners=6pt, thick, fill=gray!4,
          fit=(route)(rewrite)(sim), inner xsep=0.28cm, inner ysep=0.55cm] (runbox) {};
  \end{scope}
  \node[runlabel, above left=0.10cm and -0.05cm of runbox.north west,
        anchor=south west] {repeat for $r=1,\ldots,R$};

  \node[result, below=0.85cm of sim] (aggregate)
    {aggregate paired differences: $\delta_k^{(r)}=D_k^{(r)}-D_0^{(r)}$};

  \draw[arr] (inputs) -- (route);
  \draw[arr] (route) -- (rewrite);
  \draw[arr] (rewrite) -- (sim);
  \draw[arr] (runbox.south) -- (aggregate);
\end{tikzpicture}
\caption{Monte Carlo protocol. The demand and candidate plans are
fixed, each seed produces one baseline route realization, applies all detours
to that realization and simulates paired baseline-plan conditions.}
\label{fig:mc-workflow}
\end{figure}

\subsubsection{Performance metrics}
\label{sssec:mc_metrics}

Each \texttt{tripinfo.xml} file is parsed into a metric vector $M_c^{(r)}$ for
condition $c$ and run $r$. SUMO is configured with
\texttt{tripinfo-output.write-unfinished}, so completed trips can be separated
from unfinished or vaporized records. Aggregate performance is computed only
on completed trips. The unfinished count is kept as a diagnostic.

Time-loss delay is the primary network-efficiency metric, but it is not the
only assessment criterion. The Monte Carlo stage stores time-loss delay, mean
travel time, completed trips, unfinished trips and the number of vehicles whose
route changed after applying a plan. The evaluator then derives additional
impacted-user metrics from the route and \texttt{tripinfo.xml} files.

\begin{itemize}
  \item \textbf{Total time-loss delay} $D$ (vehicle-hours): the sum of
    SUMO \texttt{timeLoss} over completed trips, converted from seconds
    to hours,
    \begin{equation}
      D \;=\; \frac{1}{3600} \sum_{v \in \mathcal{V}_{\mathrm{done}}} \lambda_v,
      \label{eq:mc_delay}
    \end{equation}
    where $\lambda_v$ is the vehicle's SUMO \texttt{timeLoss}. This is the
    excess travel time relative to SUMO's ideal traversal of the driven route,
    so it captures both stopped waiting and moving slowdowns.

  \item \textbf{Mean travel time} $\bar{\tau}$ (minutes): the arithmetic mean
    of completed trip durations,
    \begin{equation}
      \bar{\tau} \;=\;
      \frac{1}{60\,|\mathcal{V}_{\mathrm{done}}|}
      \sum_{v \in \mathcal{V}_{\mathrm{done}}} \tau_v,
      \label{eq:mc_tt}
    \end{equation}
    where $\tau_v$ is the door-to-door duration of trip $v$ in seconds.

  \item \textbf{Completed and unfinished trips}: the number of vehicles that
    arrived within the simulation window and the number of records that did
    not. A high unfinished count flags a simulation window, demand, or
    topology problem that should be inspected before interpreting a plan.

  \item \textbf{Rerouted vehicles}: for each plan and run, the number and
    percentage of vehicles whose route sequence differs from the run's
    baseline route file. This value is computed after \texttt{duarouter}, so
    it can vary across Monte Carlo iterations.
\end{itemize}

For plan comparison, the evaluator computes paired differences such as
$\delta_k^{(r)}$ in Equation~\eqref{eq:paired_delta} run by run. These raw
paired observations are preserved in the per-run \texttt{results.json} files
and are not replaced by only top-level averages.

\subsubsection{Convergence and number of runs}
\label{sssec:mc_convergence}
The number of runs is fixed by the user. If more
precision is needed, the same output directory can be extended using a 
resume mechanism, which continues from the next missing run identifier.
